% ##############################################################################
\section{Canary bank and deterministic grading}
\label{sec:bank}

% ==============================================================================
\subsection{The candidate bank}
\label{sec:bank_content}

A \emph{canary} is a tuple $c = (\mathit{id}, v, f, p, a, \nu, \sigma)$ of a
stable identifier, a bank version $v$, a task family $f$, a prompt $p$, a
canonical answer $a$, the name of a normalizer $\nu$, and the seed $\sigma$ of
the generator that produced it. A \emph{bank} $\bank_v$ is a set of canaries of
one version. The checked-in bank \texttt{v1} contains $|\bank_1| = 300$
canaries, 50 in each of six families $\family$, stored as one JSONL file per
family. The script \texttt{generate\_canary\_bank.py} produces the bank
deterministically: record $i$ of the family with index $j \in \{0,\dots,5\}$
uses the seed $100000\,(j+1) + i$, and the generator computes the answer from
the same random parameters that it renders into the prompt. The answer is
therefore correct by construction, with no model in the loop, which supports
property P1 of Problem~\ref{prob:probing}.

Table~\ref{tab:families} describes the families. Every prompt ends with an
output-format instruction, such as ``Return only the final integer''.

\begin{table}[htbp]
  \centering
  \small
  \begin{tabularx}{\linewidth}{@{}p{0.14\linewidth}Xp{0.16\linewidth}@{}}
    \toprule
    Family & Task & Normalizer \\
    \midrule
    \texttt{arithmetic} & Evaluate a nested integer expression of the form
      $((a \cdot b) + c) \cdot (d - e) - f$ & \texttt{integer} \\
    \texttt{code-trace} & Report the integer printed by a seven-element Python
      loop with a modular condition & \texttt{integer} \\
    \texttt{modular} & Apply four updates $x \leftarrow (m x + b) \bmod q$ in
      order and report the final $x$ & \texttt{integer} \\
    \texttt{ordering} & Move one item of an eight-word list to the front,
      reverse a range, and rotate the list right & \texttt{alnum\_lower} \\
    \texttt{string} & Rotate a 14-letter string, reverse a range, and read six
      positions & \texttt{alnum\_lower} \\
    \texttt{table} & Filter six rows by an odd-$x$ and a threshold condition and
      sum $xy - p$ over the kept rows & \texttt{integer} \\
    \bottomrule
  \end{tabularx}
  \caption{Task families of the \texttt{v1} candidate bank. Each family has 50
  candidates with identifiers of the form \texttt{<family>-v1-NNN}.}
  \label{tab:families}
\end{table}

The loader \texttt{load\_canary\_bank} rejects a bank with a malformed JSON
line, an identifier outside \texttt{[a-z0-9][a-z0-9.\_-]\{2,79\}}, a duplicate
identifier, an unknown normalizer, an empty field, or more than one bank
version. The storage layer adds two further rules. It refuses to insert new
canaries into a bank version that already has a calibration run, and it
refuses to register a canary whose stored definition differs from the supplied
one. Together these rules make a bank version immutable once it has produced
evidence.

The bank is a set of \emph{candidates}. Section~\ref{sec:calibration} describes
the measurement that decides which candidates the prober may use.

% ==============================================================================
\subsection{Normalizers and grading}
\label{sec:grading}

The grader compares a response with the canonical answer after both pass
through the canary's normalizer. Table~\ref{tab:normalizers} lists the
normalizers. The \texttt{v1} bank uses \texttt{integer} for 200 canaries and
\texttt{alnum\_lower} for 100 canaries; the remaining normalizers support
future banks.

\begin{table}[htbp]
  \centering
  \small
  \begin{tabularx}{\linewidth}{@{}p{0.2\linewidth}X@{}}
    \toprule
    Normalizer & Operation \\
    \midrule
    \texttt{exact} & Trim surrounding whitespace \\
    \texttt{alnum\_lower} & Lowercase and drop every non-alphanumeric
      character \\
    \texttt{integer} & Trim whitespace, remove commas, accept exactly one
      integer with an optional sign and an optional all-zero fraction, and emit
      its canonical form; any other text is rejected \\
    \texttt{decimal} & Parse one finite decimal and emit a canonical form \\
    \texttt{json\_canonical} & Parse JSON, reject non-finite constants, and
      serialize with sorted keys and compact separators \\
    \texttt{code\_output} & Normalize line endings and trailing whitespace of
      program output \\
    \bottomrule
  \end{tabularx}
  \caption{Normalizers of version \texttt{m2-v1} of the grader.}
  \label{tab:normalizers}
\end{table}

For a canary $c$ with normalizer $\nu$ and a response $y$, the grade is
\begin{equation}
  \label{eq:grade}
  \grade(c, y) =
  \begin{cases}
    \bot & \text{if } \nu \text{ is unknown, or } \nu(a) \text{ is undefined or
      empty},\\
    0 & \text{if } y \text{ is empty, or } \nu(y) \text{ is undefined},\\
    \mathbb{1}\!\left[\nu(y) = \nu(a)\right] & \text{otherwise.}
  \end{cases}
\end{equation}
The value $\bot$ signals that the \emph{checker}, rather than the seller,
failed to produce a grade. It is never converted to $0$: the evidence rules of
Section~\ref{sec:evidence} classify a $\bot$ grade on a successful response as
a checker error, and the verdict then carries the disposition
\texttt{CHECKER\_UNHEALTHY} instead of penalizing the seller.
Table~\ref{tab:reasons} lists the reason code that accompanies each grade.

\begin{table}[htbp]
  \centering
  \small
  \begin{tabularx}{\linewidth}{@{}p{0.3\linewidth}p{0.08\linewidth}X@{}}
    \toprule
    Reason code & Grade & Condition \\
    \midrule
    \texttt{MATCH} & 1 & Normalized response equals normalized answer \\
    \texttt{VALUE\_MISMATCH} & 0 & Normalized values differ \\
    \texttt{EMPTY\_RESPONSE} & 0 & Response is absent or only whitespace \\
    \texttt{INVALID\_RESPONSE} & 0 & Response has the wrong shape for the
      normalizer, for example prose around an integer \\
    \texttt{INVALID\_EXPECTED\_ANSWER} & $\bot$ & The stored answer cannot be
      normalized or is empty \\
    \texttt{UNKNOWN\_NORMALIZER} & $\bot$ & The canary names a normalizer that
      the grader does not implement \\
    \bottomrule
  \end{tabularx}
  \caption{Reason codes returned by the grader.}
  \label{tab:reasons}
\end{table}

Two properties of this grader follow from the definition. First, grading is
strict about shape: a correct integer embedded in a sentence is graded $0$
under \texttt{integer}, so the measurement couples correctness with
instruction following. Second, \texttt{alnum\_lower} ignores whitespace, so a
space-separated word list and the same words concatenated compare equal. The
eight words of an \texttt{ordering} answer are distinct, and the effect of this
equivalence on those answers was not analyzed.

% ==============================================================================
\subsection{A pluggable measurer}
\label{sec:measurer}

The white paper treats capability measurement as an interchangeable component
(Section~\ref{sec:background_whitepaper}). The code mirrors this with a
\texttt{CapabilityMeasurer} protocol that exposes a \texttt{version} string and
a \texttt{measure} method from a canary and a response to a grade. The default
implementation, \texttt{DeterministicBenchmarkMeasurer}, delegates to
\eqref{eq:grade}. Calibration, probing, and persistence receive the measurer as
an argument and store its version with every probe request, so a verifier model
or a hosted evaluator could replace the default without changing those layers.
No such alternative measurer exists in the code base, and none is evaluated in
this paper.
